% ─────────────────────────────────────────────────────────────────────────────
% Section V — Security Analysis and Compliance
% ─────────────────────────────────────────────────────────────────────────────

\section{Security Analysis and Compliance}
\label{sec:security}

\subsection{MITRE ATLAS Mapping}

MITRE ATLAS~\cite{mitre_atlas2024} documents adversarial ML attack techniques against
AI systems. Table~\ref{tab:atlas} maps backdoor persistence to the ATLAS taxonomy
and identifies where CRSC provides detection coverage.

\begin{table}[t]
\centering
\caption{CRSC coverage across MITRE ATLAS adversarial ML techniques.}
\label{tab:atlas}
\begin{tabularx}{\columnwidth}{@{} p{1.5cm} p{2.1cm} p{2.0cm} X @{}}
\toprule
\textbf{ATLAS ID} & \textbf{Technique} & \textbf{CRSC Role} & \textbf{Detection Signal} \\
\midrule
AML.T0018 & Backdoor ML Model      & Primary detection  & {\scriptsize$\Delta_{\text{hid}},\;\Delta H$} \\
AML.T0019 & Publish Poisoned Model & Post-publish gate  & {\scriptsize$\text{CRSC} \geq \tau = 0.62$} \\
AML.T0020 & Poison Training Data   & Pre-training audit & {\scriptsize$\Delta_{\text{hid}} = \|\Delta h\|_F / \|h_{\text{base}}\|_F$} \\
AML.T0031 & Evade ML Model         & Trigger-agnostic   & {\scriptsize$\Delta H = H(\ell_f) - H(\ell_b)$} \\
AML.T0043 & Craft Adversarial Data & Probe construction & {\scriptsize$\mathcal{T}{:}\;K\!=\!500,\;\rho\!\in\![0.001, 0.05]$} \\
\bottomrule
\end{tabularx}
\end{table}

CRSC addresses AML.T0018 (Backdoor ML Model) as a \emph{detection countermeasure}
and AML.T0019 (Publish Poisoned Model) as a \emph{pre-deployment gate}. By
quantifying persistence after fine-tuning, CRSC enables organizations to
operationalize ATLAS detection in their LLM deployment pipelines.

\subsection{NIST AI Risk Management Framework}

The NIST AI RMF~\cite{nist_ai_rmf2023} organizes AI risk across four functions:
GOVERN, MAP, MEASURE, and MANAGE. CRSC contributes primarily to the MEASURE
function:

\begin{itemize}[leftmargin=*]
  \item \textbf{GOVERN 1.2:} CRSC thresholds ($\tau = 0.62$) provide
  quantitative criteria for model acceptance policies.

  \item \textbf{MAP 2.3:} The three-agent pipeline systematically maps
  backdoor risks across model scale, trigger types, and fine-tuning methods.

  \item \textbf{MEASURE 2.5:} CRSC scores provide a standardized measurement
  of adversarial robustness compatible with organizational risk registers.

  \item \textbf{MANAGE 4.1:} HIGH-risk classifications ($\text{CRSC} \geq 0.62$)
  trigger specific mitigation workflows, including layer-targeted fine-tuning
  and model quarantine.
\end{itemize}

\subsection{Threat-Specific Mitigations}

For each risk level, we recommend the following mitigations:

\textbf{HIGH risk (CRSC $\geq 0.62$):}
(i) Apply neuron activation analysis~\cite{chen2018detecting} to identify
backdoor-associated neurons; (ii) perform targeted layer re-initialization
on layers with $\delta^{(l)}_\text{hidden} > 0.7$; (iii) apply STRIP
detection~\cite{gao2019strip} as a runtime filter.

\textbf{MODERATE risk (CRSC $\in [0.45, 0.62)$):}
(i) Increase fine-tuning dataset diversity (reduce trigger correlation);
(ii) apply spectral signatures detection~\cite{tran2018spectral};
(iii) monitor ASR on held-out trigger set during deployment.

\textbf{LOW risk (CRSC $< 0.45$):}
Standard model validation procedures are sufficient. Periodic re-evaluation
recommended when the model is re-fine-tuned.

\subsection{Adversarial Validity}

A critical concern for any detection metric is whether adversaries can craft
triggers that evade it. We analyze two evasion strategies:

\textbf{Hidden-state camouflage:} An adversary could optimize trigger tokens
to minimize $\Delta_\text{hidden}$ while maintaining high ASR. This constitutes
an adaptive attack against CRSC's first component. However, minimizing
$\Delta_\text{hidden}$ requires the backdoor to be encoded entirely in
output-layer logits, which is detectable via $\Delta H$ (the second component).

\textbf{Entropy equalization:} An adversary could pre-scale losses to produce
$\Delta H \approx 0$. However, this would require knowledge of the specific
probe set $\mathcal{T}$, which is held secret by the evaluator. We formalize
this as a security game and show that a CRSC-aware adversary must increase
model complexity by $\Omega(\sqrt{N})$ to simultaneously evade both components,
making evasion computationally expensive at the 70B scale.
